\documentclass[11pt]{scrartcl}

\usepackage[stixtwo]{fontsetup}
\usepackage{amsmath}

\begin{document}
\begin{enumerate}
	\item Rendőrelvvel dolgozunk. Minden pozitív $\varepsilon$ esetén
	      elegendően nagy $n$-ekre teljesül, hogy
	      $A - \varepsilon < a_n < A + \varepsilon$, így
	      \[
		      \left(1 + \frac{A - \varepsilon}{n}\right)^n
		      < \left(1 + \frac{a_n}{n}\right)^n
		      < \left(1 + \frac{A + \varepsilon}{n}\right)^n,
	      \]
	      és limeszt véve
	      \[
		      e^{A - \varepsilon}
		      \le \liminf \left(1 + \frac{a_n}{n}\right)^n
		      \le \limsup \left(1 + \frac{a_n}{n}\right)^n
		      \le e^{A + \varepsilon}.
	      \]
	      Az $\varepsilon \to 0$ határátmenetet véve megkapjuk a bizonyítandó
	      állítást.

	\item Írjuk fel az alapot exponenciális alakban:
	      \[
		      1 + \frac{i\alpha}{n} = \rho e^{i\phi}, \quad \text{ahol} \quad
		      \rho = \sqrt{1 + \left(\frac{\alpha}{n}\right)^2}, \quad
		      \phi = \arctan\left(\frac{\alpha}{n}\right).
	      \]
	      Ezen az alakon elvégezzük a hatványozást:
	      \[
		      \left(1 + \frac{i\alpha}{n}\right)^n
		      = \rho^n e^{in\phi}
		      = \left(1 + \frac{\alpha^2 / n}{n}\right)^{\frac{n}{2}}
		      e^{in \arctan\left(\frac{\alpha}{n}\right)}
		      \xrightarrow{n \to \infty} e^{i\alpha},
	      \]
	      hiszen
	      \(
	      \left(1 + \frac{\alpha^2 / n}{n}\right)^{\frac{n}{2}} \to 1
	      \)
	      az a) állítás miatt, és
	      \(
	      e^{in \arctan\left(\frac{\alpha}{n}\right)} \to e^{i\alpha}
	      \)
	      az exponenciális függvény folytonossága és a megadott összefüggés
	      miatt.
\end{enumerate}
\end{document}
